\section*{अध्याय \ref{ch:g12:diffeq} --- अवकल समीकरण}

\begin{solution}{exo:g12:diffeq:1}
\emph{(a)} $y = C\eu^{3x}$; $y(0) = 2$ से $y = 2\eu^{3x}$ मिलता है।

\emph{(b)} $y' = -\frac12 y$, इसलिए $y = C\eu^{-x/2}$; $y(0) = -1$ से $y = -\eu^{-x/2}$ मिलता है।

\emph{(c)} साम्य $y_p = 5$; $y = C\eu^{-x} + 5$; $y(0) = 0$ से $C = -5$ मिलता है: $y = 5\left(1 - \eu^{-x}\right)$।
\end{solution}

\begin{solution}{exo:g12:diffeq:2}
$N' = kN$, इसलिए $N(t) = N_0 \eu^{kt}$। $3$ घंटों में दुगुना होने का अर्थ है $\eu^{3k} = 2$, \emph{अर्थात्}
\[
k = \frac{\ln 2}{3} \approx 0.231\ \text{h}^{-1}.
\]
\end{solution}

\begin{solution}{exo:g12:diffeq:3}
$y_p'(x) = \eu^x + x\eu^x = y_p(x) + \eu^x$: $y_p$ एक विशेष हल है (यह \cref{met:g12:diffeq:particular} की अनुनाद वाली स्थिति है)। \cref{prop:g12:diffeq:general} से हल $y = C\eu^{x} + x\eu^{x} = (C + x)\,\eu^x$,
$C \in \R$ हैं।
\end{solution}

\begin{solution}{exo:g12:diffeq:4}
$y_p = \beta\eu^x$ आज़माइए: $\beta\eu^x = -2\beta\eu^x + \eu^x$ से $3\beta = 1$ मिलता है, इसलिए $y_p = \frac13\eu^x$। व्यापक हल $y = C\eu^{-2x} + \frac13\eu^{x}$; और शर्त $y(0) = 1$
से $C = \frac23$ मिलता है:
\[
y(x) = \frac{2}{3}\,\eu^{-2x} + \frac{1}{3}\,\eu^{x}.
\]
\end{solution}

\begin{solution}{exo:g12:diffeq:5}
$\lambda = \frac{\ln 2}{5730}$ के साथ $N(t) = N_0\,\eu^{-\lambda t}$ (\cref{ex:g12:exp:halflife})। हम $\eu^{-\lambda t} = 0.6$ हल करते हैं:
\[
t = \frac{\ln(1/0.6)}{\lambda} = 5730\,\frac{\ln(5/3)}{\ln 2}
\approx 5730 \times \frac{0.5108}{0.6931} \approx 4220 \text{ वर्ष}.
\]
\end{solution}

\begin{solution}{exo:g12:diffeq:6}
\emph{1.} नमक $0.2 \times 5 = 1$ kg/min की दर से भीतर आता है। बाहर जाने वाला प्रवाह $5$
L/min पर $\frac{m}{100}$ kg/L की सांद्रता ले जाता है, \emph{अर्थात्} $\frac{m}{20}$ kg/min। अतः
$m' = 1 - \frac{m}{20}$।

\emph{2.} साम्य $m_p = 20$; $m(t) = 20 + C\eu^{-t/20}$, और $m(0) = 0$ से $C = -20$ मिलता है:
\[
m(t) = 20\left(1 - \eu^{-t/20}\right) \xrightarrow[t\to+\infty]{} 20 \text{ kg}.
\]
दीर्घ काल में टंकी की सांद्रता भीतर आते लवण-जल जितनी हो जाती है: $0.2$ kg/L
$\times$ $100$ L $= 20$ kg।
\end{solution}

\begin{solution}{exo:g12:diffeq:7}
\emph{1.} $\frac km = \frac{16}{80} = 0.2$ के साथ $v' = g - \frac{k}{m} v$। साम्य $v_\infty = \frac{mg}{k} = \frac{9.8}{0.2} = 49$ m/s; $v(t) = 49 + C\eu^{-0.2t}$, और $v(0) = 0$ से मिलता है
\[
v(t) = 49\left(1 - \eu^{-0.2 t}\right).
\]

\emph{2.} सीमांत वेग $49$ m/s ($\approx 176$ km/h)। हमें $1 - \eu^{-0.2t} = 0.95$ चाहिए, \emph{अर्थात्}
$\eu^{-0.2t} = 0.05$:
\[
t = \frac{\ln 20}{0.2} \approx \frac{3.00}{0.2} \approx 15 \text{ s}.
\]
\end{solution}

\begin{solution}{exo:g12:diffeq:8}
\emph{1.} $z = \frac1y$ से $z' = -\frac{y'}{y^2}
= -\frac{y(1-y)}{y^2} = -\frac{1-y}{y} = -\frac1y + 1 = -z + 1$ मिलता है।

\emph{2.} साम्य $z_p = 1$, इसलिए $z(t) = 1 + C\eu^{-t}$। $y(0) = \frac{1}{10}$ से $z(0) = 10$, अतः $C = 9$ और
\[
y(t) = \frac{1}{1 + 9\,\eu^{-t}} .
\]

\emph{3.} $t \to +\infty$ पर $9\eu^{-t} \to 0$ और $y(t) \to 1$। वक्र चिरपरिचित S-आकार वाला
(\emph{सिग्मॉइड}) लॉजिस्टिक वक्र है: पहले धीमी वृद्धि ($y$ छोटा, $y' \approx y$), फिर
$y = \frac12$ पर सबसे तेज़ वृद्धि (जहाँ $y' = y(1-y)$ \omterm{def:g10:functions:extrema}{अधिकतम} है), और फिर वहन-क्षमता $1$ की ओर
संतृप्ति।
\end{solution}

\begin{solution}{exo:g12:diffeq:9}
\cref{thm:g12:diffeq:affine} से $y(x) = C\eu^{ax} - \frac ba$। अतः
\[
u_{n+1} = C\eu^{a(n+1)} - \frac ba
= \eu^{a}\left(C\eu^{an} - \frac ba\right) + \frac ba\left(\eu^a - 1\right)
= q\,u_n + r
\]
जहाँ $q = \eu^{a}$ और $r = \frac{b}{a}\left(\eu^{a} - 1\right)$ हैं। शर्त $\abs q < 1$ का अर्थ है $\eu^a < 1$, \emph{अर्थात्} $a < 0$: ठीक
वही शर्त जिसके अधीन \omterm{def:g12:diffeq:ode}{अवकल समीकरण} के हल साम्य $-\frac ba$ पर \omterm{def:g12:seq:limit}{अभिसरित} होते हैं --- और
सचमुच पुनरावृत्ति का \omterm{pb:g10:functions:1}{स्थिर बिंदु} $\frac{r}{1 - q} = -\frac{b}{a}$ है।
\end{solution}

\begin{solution}{pb:g12:diffeq:1}
\textbf{1.} $y = 2\eu^{3t}$। दूसरे के लिए: साम्य $3$, इसलिए $y(0) = 0$ के साथ $y = 3 + C\eu^{-2t}$: $C = -3$:
$y = 3\left(1 - \eu^{-2t}\right)$।

\textbf{2.} $\left(y\eu^{-at}\right)' = y'\eu^{-at} -
ay\eu^{-at} = (y' - ay)\eu^{-at} = 0$: गुणनफल अचर $C$ है, इसलिए $y = C\eu^{at}$ --- कोई हल बच नहीं
निकलता।

\textbf{3.} साम्य: $y \equiv 3$। हर दूसरा हल $C \neq 0$ के साथ $3 + C\eu^{-2t}$ है: चरघातांकी पद मर
जाता है और हल सरककर $3$ पर आ जाता है --- एक स्थायी \omterm{pb:g10:functions:1}{स्थिर बिंदु}, जो \cref{pb:g11:seq:1} की
पुनरावृत्तियों के स्थिर बिंदुओं का \omterm{def:g12:limcont:continuity}{संतत} जुड़वाँ है।

\textbf{4.} $y = y_0\eu^{-t/\tau}$ तब आधा होता है जब $\eu^{-t/\tau} = \frac12$: $t = \tau\ln 2$।

\textbf{5.} सभी हल $3$ की ओर होने वाले क्षय के ऊर्ध्वाधर स्थानांतर हैं: जो
ऊपर से शुरू होते हैं वे $3$ तक गिरते हैं, जो नीचे से वे $3$ तक चढ़ते हैं,
और अचर हल $3$ वहीं बैठा रहता है --- साम्य के चारों ओर एक क़ीप।

\textbf{6.} प्रति वर्ष $\lambda = \frac{\ln 2}{5730} \approx 1.21 \times
10^{-4}$।

\textbf{7.} $\eu^{-\lambda t} = 0.2$: $t = \frac{\ln 5}{\lambda} \approx 13\,300$ वर्ष।

\textbf{8.} $t = \frac{\ln(1/0.15)}{\lambda} \approx 15\,700$ वर्ष: लास्को के साँड़ हिमयुग के अंतिम दौर के हैं।

\textbf{9.} $t = \frac{\ln(1/0.53)}{\lambda} \approx 5\,250$ वर्ष: पुरातत्त्वविदों के मान से एक जीवनकाल के भीतर ---
घड़ी काम करती है।

\textbf{10.} दस अर्धायुओं के बाद $2^{-10} \approx 0.1\,\%$ बचता है: मापन की सीमा पर। $65$ करोड़
वर्ष बाद अंश $2^{-65\,000\,000/5\,730} \approx 2^{-11\,344}$ है: किसी भी जीवाश्म में \omterm{def:g11:quad:discriminant}{मूल} भंडार का एक परमाणु भी नहीं
बचता --- डायनासोर की तारीख़ धीमी घड़ियों से बताई जाती है (पोटैशियम--आर्गन और
उसके साथी)।

\textbf{11.} $52\,\%$ से $\approx 5\,405$ वर्ष मिलते हैं और $54\,\%$ से $\approx 5\,095$: रसायन का $\pm 1\,\%$
इतिहास के मोटे तौर पर $\pm 155$ वर्ष बन जाता है --- प्रयोगशाला की यथार्थता ही
संग्रहालय की पट्टिका की यथार्थता है।

\textbf{12.} $z = T - T_a$ $z' = -kz$ को संतुष्ट करता है: $z = (T_0 - T_a)\eu^{-kt}$, और $T = T_a + z$।

\textbf{13.} आरंभिक अंतर $T_0 - T_a = 70$ है; पाँच मिनट बाद अंतर $70 - 20 = 50$ है, इसलिए $50 = 70\eu^{-5k}$:
प्रति मिनट $k = \frac{\ln(70/50)}{5} \approx 0.067$। पीने योग्य: $55 - 20 = 35 = 70\eu^{-kt}$: $t = \frac{\ln 2}{k} \approx 10.3$ मिनट।

\textbf{14.} हानि की दर अंतर $T - T_a$ के समानुपाती है: खौलती कॉफ़ी सबसे तेज़ी से
ऊष्मा बहाती है। दूध \emph{अभी} डालने से अंतर तुरंत घट जाता है, इसलिए दस मिनट
में कम ऊष्मा जाती है; अंत में डालने से कॉफ़ी पहले पूरी गति से ठंडी होती रहती
है। सबसे गर्म प्याले के लिए: दूध पहले। (अधीर मेज़बानों की ऊष्मागतिकी उल्टी
है।)

\textbf{15.} परिवेश से अंतर: आधी रात को $10$, एक घंटे बाद $8$: $\eu^{-k} = 0.8$: प्रति
घंटा $k = \ln\frac{10}{8} \approx 0.223$। मृत्यु के समय अंतर $17$ था; वह आधी रात तक घटकर $10$ हो गया:
$s = \frac{\ln(17/10)}{k} \approx 2.4$ घंटे। मृत्यु का समय: लगभग $21{:}40$।

\textbf{16.} गरम या हवादार कमरा ($T_a$ अचर नहीं) घड़ी को किसी भी दिशा में मोड़
देता है; कपड़े या शरीर का द्रव्यमान $k$ बदल देते हैं (मृत्यु-परीक्षक की
सारणियाँ इसका समायोजन करती हैं) --- और ग़लत $k$ पूरे \omterm{def:g10:numbers:interval}{अंतराल} को खींच देता
है; और कहीं और से लाया गया शव क्षय के बीच में $T_a$ को फिर से शून्य कर देता
है, जिससे मृत्यु पहले या बाद की लगने लगती है। \omterm{def:g10:algebra:equation}{समीकरण} ईमानदार है; जोखिम
मान्यताएँ उठाती हैं।

\textbf{17.} $z = \frac1y$ $z' = 1 - z$ मानता है: $z = 1 + 9\eu^{-t}$ ($z(0) = 10$ से), इसलिए $y = \frac{1}{1 + 9\eu^{-t}}$।
नति-परिवर्तन $y = \frac12$ पर: $9\eu^{-t} = 1$: $t = \ln 9 \approx 2.2$ --- महामारी के मुड़ने की तारीख़, सीधे
सूत्र से।

\textbf{18.} सीमांत वेग: $v = 20$ m/s पर $v' = 0$। हल: $v(t) = 20\left(1 - \eu^{-t/2}\right)$। तब $0.95$: $\eu^{-t/2} = 0.05$:
$t = 2\ln 20 \approx 6$ s --- वर्षा की बूँदें सेकंडों में अपनी क्रूज़-चाल पा लेती हैं, और इसीलिए
वर्षा जान नहीं लेती।

\textbf{19.} स्थायी अवस्था: $c = 8$। उसका आधा: $4 = 8\left(1 - \eu^{-t/2}\right)$: $\eu^{-t/2} = \frac12$: $t = 2\ln 2 \approx 1.4$ घंटे। ड्रिप
ऋण का \omterm{def:g12:limcont:continuity}{संतत} जुड़वाँ है: अचर आवक, समानुपाती जावक --- और \cref{exo:g12:diffeq:9} शब्दकोश को
ठीक-ठीक बना देता है।

\textbf{20.} क्षय ($a < 0$, $b = 0$), शीतलन ($y = T - T_a$), गिरना ($a < 0$, $b > 0$: नीचे से
साम्य की ओर पहुँच), ख़ुराक (वही, पर शिरा में): एक \omterm{def:g10:algebra:equation}{समीकरण}, चार संसार।
प्रक्रिया-चक्र: परिवर्तन का नियम बताइए; \omterm{def:g10:algebra:equation}{समीकरण} लिखिए; हल कीजिए (\cref{met:g12:diffeq:solve}); मापों
पर अचर अंशांकित कीजिए; भविष्यवाणी कीजिए; और फिर मान्यताओं की जाँच कीजिए (प्रश्न
16)। दोलन के लिए $y'' = -\omega^2 y$ चाहिए --- जिससे भेंट, जिसका हल और जिसका झूला \cref{pb:g12:trigo:1} में
है।
\end{solution}
